\chapter{Training Configurations and Hyperparameters}
\label{ch:appendix_b}

This appendix reports the complete training configurations for the two supervised vision models in \cref{sec:vision_inference_strategies} and the \acs{LightGBM} classifiers in \cref{sec:downstream_model}.
All values were obtained from the training records, including the history JSON file for each fold and the run arguments stored in the checkpoints, rather than inferred from software defaults.
All model training and \ac{VLM} inference used one NVIDIA GeForce RTX 4060 Laptop graphics processing unit (GPU) with 8\,GB of video memory (VRAM).

\section{Supervised Vision Configurations}
\label{sec:appendix_b_vision_configs}

The two supervised configurations use the same training procedure: the AdamW optimiser; weight decay of 0.01, excluding normalisation parameters and biases; a cosine-annealing learning-rate schedule with $T_{\max} = 30$; a maximum of 30 epochs; early stopping based on validation building type macro F1 with a patience of 10 epochs; and fp16 automatic mixed precision (AMP).
The shared trunks use Gaussian error linear unit (GELU) activation.
Input images are resized to 224 $\times$ 224 pixels and normalised using ImageNet statistics (mean = (0.485, 0.456, 0.406); standard deviation = (0.229, 0.224, 0.225)).
Training augmentation comprises \texttt{RandomResizedCrop} with a scale range of 0.7--1.0, horizontal flipping with $p = 0.5$, and \texttt{ColorJitter} with brightness, contrast, and saturation set to 0.2.
Validation and holdout inference use deterministic resizing to a short side of 256 pixels followed by a 224-pixel centre crop.
Building type cross-entropy is weighted inversely to the frequencies of SFH, TH, MFH, and AB within each training fold.
Construction year and floor count are standardised using the mean and standard deviation of the training fold before mean squared error is calculated.

\begin{table}[htbp]
	\centering
	\caption{Training hyperparameters of the two supervised vision configurations.}
	\label{tab:vision_hyperparams}
	\footnotesize
	\begin{tblr}{
			width = \linewidth,
			colspec = {Q[l,m,3.8cm] X[l] X[l]},
			colsep = 4pt,
			row{1} = {font=\bfseries},
			rowsep = 2pt,
		}
		\toprule
		Setting                          & ResNet-50 full fine-tuning                                            & Frozen DINOv2 \acs{ViT}-B/14                                                     \\
		\midrule
		Pre-trained weights (\texttt{timm}) & \texttt{resnet50} (ImageNet)                                       & \texttt{vit\_base\_\allowbreak patch14\_\allowbreak dinov2.\allowbreak lvd142m}   \\
		Backbone feature dimension       & 2048                                                                  & 768                                                                              \\
		Shared trunk                     & Linear (2048 $\rightarrow$ 256) + GELU + Dropout 0.3                  & Linear (768 $\rightarrow$ 256) + GELU + Dropout 0.3                              \\
		Parameters updated               & Backbone, shared trunk, and output heads                              & Shared trunk and output heads only                                               \\
		Total parameters                 & 24.03\,M                                                              & 85.92\,M                                                                         \\
		Trainable parameters             & 24.03\,M (100\%)                                                      & 0.20\,M (approximately 0.2\%)                                                    \\
		Learning rate                    & Backbone $1\times10^{-4}$; trunk/heads $5\times10^{-4}$ (grouped)     & $5\times10^{-4}$                                                                 \\
		Batch size (buildings)           & 8                                                                     & 32                                                                               \\
		Optimizer / weight decay         & AdamW / 0.01                                                          & AdamW / 0.01                                                                     \\
		Learning-rate schedule           & Cosine annealing ($T_{\max} = 30$)                                    & Cosine annealing ($T_{\max} = 30$)                                               \\
		Max epochs / early-stop patience & 30 / 10                                                               & 30 / 10                                                                          \\
		Mixed precision                  & fp16 AMP                                                              & fp16 AMP                                                                         \\
		Padding handling                 & Valid images only enter the backbone to preserve BatchNorm statistics & Zero-padding with masked pooling                                                 \\
		\bottomrule
	\end{tblr}
\end{table}

The M2 direct energy class configurations follow the same training procedure as \cref{tab:vision_hyperparams}, except that the three attribute heads are replaced by one energy class classification head and early stopping uses validation energy class macro F1.
M2-ResNet-50 is fine-tuned end to end using the learning rates and batch size reported above.
Because the DINOv2 backbone is frozen, M2-DINOv2 first extracts and caches building-level mean features using deterministic resizing and centre cropping without augmentation.
The same classification-network structure is then trained on the cached features using AdamW, a learning rate of $5\times10^{-4}$, weight decay of 0.01, cosine annealing, early-stopping patience of 10 epochs, and a minibatch size of 256.
Unlike the supervised attribute-extraction procedure in \cref{tab:vision_hyperparams}, this cached-feature route does not use training-time image augmentation.

\section{LightGBM Configuration}
\label{sec:appendix_b_lightgbm}

\begin{table}[htbp]
	\centering
	\caption{Fixed LightGBM hyperparameters shared by all downstream runs.}
	\label{tab:lightgbm_hyperparams}
	\small
	\begin{tblr}{
			width = \linewidth,
			colspec = {Q[l,m,5.0cm] X[l]},
			colsep = 4pt,
			row{1} = {font=\bfseries},
			rowsep = 2pt,
		}
		\toprule
		Hyperparameter              & Value                                            \\
		\midrule
		objective                   & \texttt{multiclass} (seven-class) or \texttt{binary} \\
		\texttt{n\_estimators}      & 400                                              \\
		\texttt{learning\_rate}     & 0.05                                             \\
		\texttt{num\_leaves}        & 31                                               \\
		\texttt{max\_depth}         & $-1$ (unlimited)                                 \\
		\texttt{min\_child\_samples} & 20                                              \\
		\texttt{subsample}          & 0.8 per iteration                                \\
		\texttt{colsample\_bytree}  & 0.8                                              \\
		\texttt{reg\_lambda} (L2)   & 1.0                                              \\
		\texttt{class\_weight}      & None                                             \\
		\texttt{random\_state}      & 42                                               \\
		\bottomrule
	\end{tblr}
\end{table}

The same hyperparameters are used for all attribute-based routes and both energy class tasks.
Only the objective changes with the task; no route-specific tuning is performed.
